\chapter{Álgebras cuadráticas de orden \(9\)}
\label{chap:algebras-cuadraticas-orden-nueve}
\index{álgebra cuadrática}
\index{cuerpo finito \(\mathbb F_9\)}

La identidad \(A_W^2=-I_6\), obtenida en la sección anterior por reducción
de la matriz de conferencia de Paley, contiene más información que una
ortogonalidad modular. Determina una extensión de escalares de grado dos,
convierte el espacio ternario de seis dimensiones en un espacio de dimensión
tres sobre el cuerpo de nueve elementos y permite escribir el código de
Golay ternario como el grafo de una unidad cuadrática. Esta sección desarrolla
esas consecuencias y las distingue de la aritmética de característica nueve
empleada por la rama de Hensel.

\section{Estructura de \texorpdfstring{\(\mathbb F_9\)}{F9} inducida por Paley}

El polinomio \(X^2+1\) no tiene raíces en \(\mathbb F_3\): sus valores en
\(0,1,2\) son \(1,2,2\). Es, por tanto, irreducible, y el cociente
\[
 \mathbb F_9
 :=\mathbb F_3[\iota]/(\iota^2+1)
\]
es un cuerpo de nueve elementos.

\begin{theorem}[Estructura cuadrática inducida por \(A_W\)]
\label{thm:f9-tres-desde-paley}
Sea \(V=\mathbb F_3^6\), escrito mediante vectores fila, y sea
\(J(v)=vA_W\). La regla
\[
 (a+b\iota)v:=av+bJ(v),
 \qquad a,b\in\mathbb F_3,
\]
convierte \(V\) en un espacio vectorial de dimensión tres sobre
\(\mathbb F_9\). En las coordenadas fijadas por la construcción de Paley,
\[
 \mathbb F_3^6\cong\mathbb F_9^3.
\]
\end{theorem}

\begin{proof}
La igualdad \(J^2=-I\) hace que la evaluación
\(\mathbb F_3[X]\to\operatorname{End}_{\mathbb F_3}(V)\), \(X\mapsto J\),
se anule sobre el ideal \((X^2+1)\). La acción factoriza, por ello, a través
del cuerpo \(\mathbb F_9\). Como
\([\,\mathbb F_9:\mathbb F_3\,]=2\) y
\(\dim_{\mathbb F_3}V=6\), se obtiene
\(\dim_{\mathbb F_9}V=3\).
\end{proof}

El isomorfismo depende del sistema de coordenadas que fija \(A_W\). La
afirmación no se deduce del mero cardinal \(3^6=9^3\): la matriz exhibe la
multiplicación por todos los escalares \(a+b\iota\) y proporciona, por
tanto, la estructura lineal completa.

\section{El código ternario como grafo isotrópico}

La aplicación generadora del código es
\[
 c(q)=(q,qA_W).
\]
Después de identificar \(J\) con la multiplicación por \(\iota\), su imagen
admite la escritura
\[
 C_W=\{(q,\iota q):q\in\mathbb F_9^3\},
\]
donde ambas componentes se consideran, por restricción de escalares, como
vectores de \(\mathbb F_3^6\).

\begin{proposition}[Maximalidad isotrópica]
\label{prop:golay-grafo-isotropico}
Respecto de la forma bilineal estándar de \(\mathbb F_3^{12}\), el
subespacio \(C_W\) es totalmente isótropo y maximal entre los subespacios
totalmente isótropos.
\end{proposition}

\begin{proof}
La matriz generadora \(G=(I_6\mid A_W)\) satisface
\[
 GG^{\mathsf T}=I_6+A_WA_W^{\mathsf T}=I_6+2I_6=0.
\]
El espacio generado es, por consiguiente, totalmente isótropo. Tiene
dimensión seis, la mitad de la dimensión del espacio ambiente no degenerado,
y es, por tanto, maximal; de hecho, coincide con su dual ortogonal.
\end{proof}

La raíz de \(-I\) aparece así como una transformación finita de incidencia.
No se introduce el cuerpo complejo analítico en el sistema: se obtiene una
estructura cuadrática sobre \(\mathbb F_3\), suficiente para organizar la
polaridad, el código y sus dos espacios propios después de extender
escalares.

\section{Característica \(9\) y característica \(3\)}

Dos conjuntos de \(729\) estados intervienen en la construcción:
\[
 (\mathbb Z/9\mathbb Z)^3
 \qquad\text{y}\qquad
 \mathbb F_9^3.
\]
No son isomorfos como grupos aditivos ni como anillos. El primero tiene
característica nueve y contiene elementos de orden nueve; el segundo tiene
característica tres y todo su grupo aditivo tiene exponente tres.

Para \(a,b\in\mathbb F_3\), eligiendo representantes
\(\widetilde a,\widetilde b\in\{0,1,2\}\), sea
\[
 \beta(a,b)=\widetilde a+3\widetilde b\pmod 9.
\]
La aplicación coordenada inducida por \(\beta\) es una biyección de
conjuntos entre \(\mathbb F_3^6\) y \((\mathbb Z/9\mathbb Z)^3\), pero no
respeta la suma. En efecto, si
\[
 \kappa(a,c)
 :=\left\lfloor
 \frac{\widetilde a+\widetilde c}{3}
 \right\rfloor,
\]
entonces
\[
 \beta(a,b)+\beta(c,d)
 =\beta\bigl(a+c,\ b+d+\kappa(a,c)\bigr),
\]
con las coordenadas del segundo miembro reducidas módulo tres. Por ejemplo,
\[
 \beta(2,0)+\beta(2,0)=\beta(1,1),
\]
mientras que la suma coordenada en \(\mathbb F_3^2\) daría \((1,0)\). La
segunda coordenada de \(\beta(1,1)\) conserva el cociente de la división
euclídea; este término es precisamente el dato que desaparece bajo la
reducción componente a componente.

Esta distinción tipa las dos lecturas del cardinal \(729\). La realización
\((\mathbb Z/9\mathbb Z)^3\) conserva residuo y cociente; la realización
\(\mathbb F_9^3\) convierte cada par ternario en un escalar de una extensión
cuadrática. Una prolonga la aritmética con memoria de división; la otra
organiza la incidencia excepcional.

\begin{proposition}[Dos lecturas tipadas del mismo soporte ternario]
\label{prop:dos-lecturas-tipadas-729}
Fijadas la sección conjuntista empleada en \(\beta\), la base, la
polarización y el operador \(J\) del
\cref{thm:f9-tres-desde-paley}, existe el diagrama de biyecciones de conjuntos
\[
\boxed{
(\mathbb Z/9\mathbb Z)^3
\xleftarrow{\;\beta^{\times3}\;}
\mathbb F_3^6
\xrightarrow{\;\iota_J\;}
\mathbb F_9^3.}
\]
La flecha izquierda transporta la suma nonádica torcida por el cociclo de
acarreo \(\kappa\); la flecha derecha transporta la estructura lineal sobre
\(\mathbb F_9\) determinada por \(J^2=-I\). Las dos flechas comparten el
soporte coordenado \(\mathbb F_3^6\), pero no identifican sus leyes
algebraicas.
\end{proposition}

\begin{proof}
La fórmula
\[
 \beta(a,b)+\beta(c,d)
 =\beta\bigl(a+c,b+d+\kappa(a,c)\bigr)
\]
prueba que \(\beta^{\times3}\) es una biyección con suma torcida y que no es
un isomorfismo desde la suma componente a componente de
\(\mathbb F_3^6\). Por el \cref{thm:f9-tres-desde-paley}, la elección de
\(J\) convierte ese mismo espacio ternario de dimensión seis en un espacio
de dimensión tres sobre
\(\mathbb F_3[J]\cong\mathbb F_9\), lo que define \(\iota_J\).
\end{proof}

La proposición es una síntesis demostrada con estructura de partida
explícita. No afirma una identificación canónica entre
\((\mathbb Z/9\mathbb Z)^3\) y \(\mathbb F_9^3\): el primero contiene
elementos aditivos de orden nueve, mientras que el segundo posee exponente
aditivo tres. El cociclo mide precisamente la información que se perdería al
confundir ambas leyes.

\section{Clasificación de las álgebras cuadráticas}

Sobre un cuerpo de característica distinta de dos, un álgebra cuadrática
con unidad se clasifica por el tipo de su polinomio definitorio. Sobre
\(\mathbb F_3\) pueden elegirse los tres representantes
\[
 \mathcal Q_{\mathrm e}
 =\mathbb F_3[t]/(t^2+1)\cong\mathbb F_9,
\]
\[
 \mathcal Q_{\mathrm n}
 =\mathbb F_3[\varepsilon]/(\varepsilon^2),
 \qquad
 \mathcal Q_{\mathrm s}
 =\mathbb F_3[r]/(r^2-1)
 \cong\mathbb F_3\times\mathbb F_3.
\]

\begin{theorem}[Clasificación de las álgebras cuadráticas sobre \(\mathbb F_3\)]
\label{thm:clasificacion-cuadratica-nonadica}
Las relaciones
\[
 J^2=-I,\qquad N^2=0,\qquad R^2=I
\]
definen acciones de
\(\mathcal Q_{\mathrm e}\), \(\mathcal Q_{\mathrm n}\) y
\(\mathcal Q_{\mathrm s}\), respectivamente. Estas tres álgebras poseen
nueve elementos y son mutuamente no isomorfas: la primera es un cuerpo; la
segunda contiene nilpotentes no nulos; la tercera contiene idempotentes no
triviales.
\end{theorem}

\begin{proof}
Cada relación hace factorizar la evaluación de
\(\mathbb F_3[t]\) por el ideal correspondiente. La irreducibilidad de
\(t^2+1\) prueba que \(\mathcal Q_{\mathrm e}\) es un cuerpo. En
\(\mathcal Q_{\mathrm n}\), la clase \(\varepsilon\) es no nula y satisface
\(\varepsilon^2=0\). En \(\mathcal Q_{\mathrm s}\), como \(2^{-1}=2\) en
\(\mathbb F_3\), los elementos
\[
 e_+=2(1+r),\qquad e_-=2(1-r)
\]
son idempotentes ortogonales no triviales y suman la unidad. Ser un cuerpo,
poseer un nilpotente no nulo y poseer idempotentes no triviales son
propiedades invariantes por isomorfismo.
\end{proof}

En tres copias de la representación regular de dimensión dos se obtienen
tres acciones de dimensión seis sobre \(\mathbb F_3\). La relación
\(J^2=-I\) corresponde a la polaridad de Paley; la relación \(N^2=0\)
describe la linealización de una extensión con cociente; la relación
\(R^2=I\) separa dos sumandos mediante los idempotentes \(e_\pm\). Esta
clasificación proporciona el marco algebraico finito en el que las
estructuras elíptica, no reducida y escindida se distinguen sin apelar a una
geometría continua.

La construcción no reemplaza a \(\mathbb C\): el cuerpo \(\mathbb F_9\) no
posee la topología, la completitud ni la estructura analítica compleja. El
resultado exacto es otro: una raíz cuadrada de \(-I\) y la correspondiente
extensión de escalares emergen de la matriz de incidencia ya construida
dentro del sistema ternario.

\section{Fuente cuadrática integral y especializaciones finita y real}
\label{sec:fuente-cuadratica-integral-rev3}

La raíz finita y la raíz operatorial real se coordinan mediante una
presentación integral común. Sea
\begin{equation}
 \mathcal Q_{\mathbb Z}
 :=\mathbb Z[u]/(u^2+1).
 \label{eq:fuente-cuadratica-integral-rev3}
\end{equation}
Sus cambios de base son
\[
 \mathcal Q_{\mathbb Z}\otimes_{\mathbb Z}\mathbb F_3
 \cong\mathbb F_3[u]/(u^2+1)
 \cong\mathbb F_9,
\]
y
\[
 \mathcal Q_{\mathbb Z}\otimes_{\mathbb Z}\mathbb R
 \cong\mathbb R[u]/(u^2+1).
\]

\begin{theorem}[Dos realizaciones de una presentación cuadrática integral]
\label{thm:dos-realizaciones-fuente-cuadratica-rev3}
La representación finita \(u\mapsto A_W\) sobre
\(\mathbb F_3^6\) y la representación real \(u\mapsto J_{\mathrm e}\) sobre
\(V_{\mathrm e}\), construida en el
\cref{thm:tres-generadores-concretos-trit}, realizan, después de sus
respectivos cambios de base, la misma relación integral
\[
 u^2+1=0.
\]
\end{theorem}

\begin{proof}
La igualdad \(A_W^2+I_6=0\) hace factorizar la evaluación
\(\mathbb F_3[u]\to\operatorname{End}_{\mathbb F_3}(\mathbb F_3^6)\) por
\((u^2+1)\). Del mismo modo, \(J_{\mathrm e}^2+I=0\) hace factorizar
\(\mathbb R[u]\to\operatorname{End}_{\mathbb R}(V_{\mathrm e})\). Las dos
representaciones proceden de especializaciones de
\(\mathcal Q_{\mathbb Z}\); no existe ni se utiliza un homomorfismo unital
\(\mathbb F_9\to\mathbb R\).
\end{proof}

La fuente común coordina relaciones, no identifica módulos de
características distintas. La rama finita conserva incidencia; la rama real
incorpora la estructura analítica necesaria para la exponencial operatorial.

\section{\(3\) generadores concretos de los regímenes TRIT}
\label{sec:tres-generadores-concretos-trit}
\index{TRIT!generadores concretos}
\index{generador elíptico}
\index{generador parabólico}
\index{generador hiperbólico}

El \cref{chap:trit-algebras} fija los tres tipos cuadráticos abstractos
\(J^2=-I\), \(J^2=0\) y \(J^2=I\), mientras el
\cref{thm:clasificacion-cuadratica-nonadica} proporciona sus modelos finitos
sobre \(\mathbb F_3\). Esos objetos no deben confundirse con las
realizaciones reales siguientes. Aquí se fijan tres espacios reales y tres
operadores concretos, cada uno obtenido de su antecedente causal declarado.

Sea
\[
V_{\mathrm e}=A_2\otimes_{\mathbb Z}\mathbb R,
\qquad
V_{\mathrm p}=\mathbb R^2,
\qquad
V_{\mathrm h}=\mathbb R^2.
\]
Sobre \(V_{\mathrm e}\) actúa el Coxeter \(C\) de la
\cref{prop:coxeter-a2-primario}. Sobre \(V_{\mathrm p}\) se fija
\[
N=
\begin{pmatrix}
0&1\\
0&0
\end{pmatrix}.
\]
Finalmente, el \cref{thm:descenso-necesario-fav} construye
\[
F_{\rm av}=
\begin{pmatrix}
0&1\\
1&1
\end{pmatrix},
\qquad
A:=F_{\rm av}^{\,2}=
\begin{pmatrix}
1&1\\
1&2
\end{pmatrix}
\]
sobre \(V_{\mathrm h}\).

\begin{theorem}[Tres realizaciones concretas de los regímenes TRIT]
\label{thm:tres-generadores-concretos-trit}
Defínanse
\[
J_{\mathrm e}:=\frac{2C+I}{\sqrt3},
\qquad
J_{\mathrm p}:=N,
\qquad
J_{\mathrm h}:=\frac{2A-3I}{\sqrt5}.
\]
Entonces
\[
J_{\mathrm e}^{\,2}=-I_{V_{\mathrm e}},
\qquad
J_{\mathrm p}^{\,2}=0,
\qquad
J_{\mathrm h}^{\,2}=I_{V_{\mathrm h}}.
\]
Por consiguiente, para \(t\in\mathbb R\),
\[
\exp(tJ_{\mathrm e})
=\cos t\,I+\sin t\,J_{\mathrm e},
\qquad
\exp(tJ_{\mathrm p})
=I+tJ_{\mathrm p},
\]
\[
\exp(tJ_{\mathrm h})
=\cosh t\,I+\sinh t\,J_{\mathrm h}.
\]

Además, en la rama elíptica,
\[
C=-\frac12I+\frac{\sqrt3}{2}J_{\mathrm e},
\]
y, para \(n\in\mathbb Z\),
\[
C^n=
\cos\!\left(\frac{2\pi n}{3}\right)I
+\sin\!\left(\frac{2\pi n}{3}\right)J_{\mathrm e}.
\]
En la rama hiperbólica, la forma
\[
G=
\begin{pmatrix}
2&1\\
1&-2
\end{pmatrix}
\]
tiene firma \((1,1)\) y satisface
\[
A^{\mathsf T}GA=G,\qquad
\det A=1,\qquad
\operatorname{tr}A=3.
\]
Si
\[
\eta:=\operatorname{arcosh}\frac32=2\log\varphi,
\]
entonces
\[
A=\cosh(\eta)I+\sinh(\eta)J_{\mathrm h},
\qquad
A^n=\cosh(n\eta)I+\sinh(n\eta)J_{\mathrm h}.
\]
\end{theorem}

\begin{proof}
El polinomio característico del Coxeter es \(X^2+X+1\); por
Cayley--Hamilton,
\[
C^2+C+I=0.
\]
De aquí,
\[
(2C+I)^2=4C^2+4C+I=-3I,
\]
y, por tanto, \(J_{\mathrm e}^2=-I\). La identidad para \(C^n\) sigue de
\(C=-\frac12I+\frac{\sqrt3}{2}J_{\mathrm e}\).

La multiplicación directa da \(N^2=0\). Por ello, la serie exponencial se
trunca exactamente:
\[
\exp(tN)=I+tN.
\]

Para la tercera rama,
\[
A^{\mathsf T}GA=G,\qquad
\det G=-5<0,
\]
de modo que \(G\) es no degenerada y posee firma \((1,1)\). Asimismo,
\[
(2A-3I)^2=5I,
\]
por lo que \(J_{\mathrm h}^2=I\). Como
\[
\cosh(2\log\varphi)=\frac32,
\qquad
\sinh(2\log\varphi)=\frac{\sqrt5}{2},
\]
se obtiene
\[
A=\frac32I+\frac{\sqrt5}{2}J_{\mathrm h}
=\cosh(\eta)I+\sinh(\eta)J_{\mathrm h}.
\]
La fórmula para \(A^n\) resulta de la ley exponencial. Las tres expresiones
generales se obtienen, en cada caso, separando en la serie exponencial las
potencias pares e impares.
\end{proof}

\begin{corollary}[Ecuación de Euler extendida por el parámetro trítico]
\label{cor:euler-extendida-trit-fibonacci}
Para \(\tau\in\{+1,0,-1\}\), sea \(J_\tau\) uno de los generadores
anteriores, normalizado por
\[
 J_\tau^2=-\tau I.
\]
Defínanse las funciones de la familia cuadrática
\[
 C_\tau(t)=
 \begin{cases}
  \cos t,&\tau=+1,\\
  1,&\tau=0,\\
  \cosh t,&\tau=-1,
 \end{cases}
 \qquad
 S_\tau(t)=
 \begin{cases}
  \sin t,&\tau=+1,\\
  t,&\tau=0,\\
  \sinh t,&\tau=-1.
 \end{cases}
\]
Entonces las tres ramas satisfacen una única identidad exponencial:
\begin{equation}
 \boxed{\exp(tJ_\tau)=C_\tau(t)I+S_\tau(t)J_\tau.}
 \label{eq:euler-extendida-trit-fibonacci}
\end{equation}
En particular,
\[
 \exp(\pi J_{+1})+I=0,
 \qquad
 \exp(J_0)=I+J_0,
 \qquad
 \exp\!\bigl(2\log\varphi\,J_{-1}\bigr)=F_{\rm av}^{\,2}.
\]
La primera igualdad es el cierre matricial de Euler; la segunda realiza el
transporte parabólico de umbral; la tercera incorpora la autoescala de
Fibonacci mediante los autovalores \(\varphi\) y \(-\varphi^{-1}\).
\end{corollary}

\begin{proof}
La identidad \eqref{eq:euler-extendida-trit-fibonacci} se obtiene separando
las potencias pares e impares de la serie exponencial y usando
\(J_\tau^2=-\tau I\). Las tres especializaciones siguen de
\(J_{+1}^2=-I\), \(J_0^2=0\) y de la igualdad ya demostrada
\(F_{\rm av}^{\,2}=\exp(2\log\varphi\,J_{-1})\).
\end{proof}

\paragraph{\(3\) realizaciones exactas de los regímenes cuadráticos del TRIT}
Las ramas elíptica, parabólica e hiperbólica realizan, respectivamente, el
Coxeter interno de \(A_2\), el generador nilpotente y la incidencia
\(F_{\rm av}\) inducida por el calendario TPK. Las identidades matriciales
son exactas en los espacios y calibres declarados.

\paragraph{Condiciones de realización de los \(3\) regímenes cuadráticos del TRIT y obstrucciones de no degeneración algebraica}
El tipo TRIT aislado organiza las tres posibilidades cuadráticas, pero no
selecciona por sí solo los operadores \(C\), \(N\) y \(F_{\rm av}\): la
selección utiliza, respectivamente, el Coxeter de APP, el modelo nilpotente
de transporte y el descenso de modos de TPK. La aparición de \(\pi\) en la
fórmula espectral elíptica pertenece a la escritura analítica del ciclo y
no sustituye la generación regional de \(\pi_{\rm HMT}\). La identidad
\(A^{\mathsf T}GA=G\) construye una transformación de Lorentz exacta en
dimensión \(1+1\); no identifica por sí sola esta pantalla con el
espacio-tiempo físico \(3+1\). Un fallo de
\[
(2C+I)^2=-3I,\qquad N^2=0,\qquad
(2A-3I)^2=5I,\qquad A^{\mathsf T}GA=G
\]
refutaría la realización correspondiente.
